% Requires: \usepackage{booktabs} and \usepackage{amsmath}
% Numbers: attentive-pooling probe, epoch 20.
% runs/channels/probe_attentive, runs/ablation/probe_attentive.

\subsection{What the reconstruction target must contain}
\label{sec:ablation-channels}

Our representation gives each degree of freedom five channels per frame: the joint
angle encoded as $(\sin\theta, \cos\theta)$, its first and second time derivatives
$(\dot\theta, \ddot\theta)$, and the generalized joint torque $\tau$ recovered by
inverse dynamics. The masked-reconstruction objective need not predict all of them.
We therefore ablate the \emph{loss} channels while leaving the \emph{input} untouched:
the encoder always reads all five channels, and only the decoder head and the
reconstruction loss are narrowed. Each arm is pretrained identically and evaluated by
the frozen-backbone linear probe of Section~\ref{sec:probe}.

The comparison is not simply kinematic against dynamic. The equations of motion relate
the two directly,
\begin{equation}
\tau = M(\theta)\,\ddot\theta + C(\theta,\dot\theta) + G(\theta),
\label{eq:eom}
\end{equation}
so acceleration is the kinematic half of the dynamics and a partial proxy for torque.
Two arms are included specifically to control for this: $\{\sin\theta,\ddot\theta\}$,
which pits torque against its proxy at matched head width, and
$\{\sin\theta,\cos\theta,\dot\theta,\ddot\theta\}$, which is every channel a
kinematics-only method could compute and therefore the setting our objective must beat.

\begin{table}[t]
\centering
\small
\begin{tabular}{llcccc}
\toprule
Loss channels & Content & $d_\text{head}$ & macro mAP & micro mAP & macro F1 \\
\midrule
\multicolumn{2}{l}{\emph{chance (label prevalence)}}   & --- & 0.042 & ---   & ---   \\
\multicolumn{2}{l}{\emph{random frozen encoder}}       & --- & 0.128 & 0.334 & 0.021 \\
\midrule
$\tau$                                        & dynamic          &  8 & 0.153 & 0.367 & 0.041 \\
$\sin\theta,\cos\theta$                       & kinematic        & 16 & 0.180 & 0.392 & 0.036 \\
$\sin\theta,\ddot\theta$                      & kin.\ + proxy    & 16 & 0.238 & 0.438 & 0.085 \\
$\sin\theta,\tau$                             & kin.\ + dynamic  & 16 & 0.259 & 0.458 & 0.100 \\
$\sin\theta,\cos\theta,\tau$                  & kin.\ + dynamic  & 24 & 0.261 & 0.460 & 0.116 \\
$\sin\theta,\cos\theta,\dot\theta,\ddot\theta$& kinematic only   & 32 & 0.216 & 0.426 & 0.074 \\
all five                                      & kin.\ + dynamic  & 40 & \textbf{0.297} & \textbf{0.493} & \textbf{0.143} \\
\bottomrule
\end{tabular}
\caption{Linear-probe transfer as a function of which channels the pretext loss scores,
under attentive pooling. $d_\text{head}$ is the width of the decoder's prediction head,
i.e.\ the number of values regressed per masked token. The all-five row is the mean of
two seeds; the seed-to-seed spread is $0.008$ macro mAP, which we take as the noise
floor. Absolute values depend on the pooling head; the ordering does not
(Section~\ref{sec:pooling}).}
\label{tab:channels}
\end{table}

\paragraph{Neither signal alone transfers well.} Reconstructing the joint angle and
nothing else reaches $0.180$ macro mAP, recovering $31\%$ of the gap between the random
encoder and the full objective; torque alone recovers $15\%$. Both fall far short in
isolation.

\paragraph{Together they exceed the sum of their parts.} Adding $\tau$ to the angle
raises macro mAP from $0.180$ to $0.259$ and macro F1 from $0.036$ to $0.100$, ten times
the seed noise floor. Separately the two account for $31\%$ and $15\%$ of the macro mAP
gap; jointly for $78\%$. On macro F1 the effect is starker still, $12\%$ and $17\%$
separately against $65\%$ together, and it holds on every metric we track. Kinematics
and dynamics are therefore complementary rather than redundant, and their
complementarity is most pronounced on the rare classes that macro-averaged metrics
weight equally --- consistent with the observation that the same joint trajectory
executed under different loads constitutes a different action, a distinction invisible
to a purely kinematic target.

\paragraph{Kinematics alone saturates well below the full objective.} The strongest
purely kinematic arm, scoring the angle together with both its derivatives, reaches
$0.216$ --- only $53\%$ of the gap, despite a $32$-dimensional head. This is the arm
that corresponds to what a method without inverse dynamics can supervise, and it is
beaten by $\{\sin\theta,\tau\}$, which reaches $0.259$ with half the head width and a
quarter of the channels. Access to torque is thus not a marginal refinement of a
kinematic objective; it is worth more than every derivative of the trajectory combined.

\paragraph{Acceleration proxies torque, but does not replace it.} Equation~\ref{eq:eom}
predicts that $\ddot\theta$ should carry part of the dynamic signal, and it does:
$\{\sin\theta,\ddot\theta\}$ reaches $0.238$ against the angle-only $0.180$, recovering
about $70\%$ of what torque contributes at the same head width. The remaining gap to
$\{\sin\theta,\tau\}$ is $0.021$, or $2.7$ times the noise floor. Acceleration is
therefore a genuine but incomplete substitute, which is what the mass, Coriolis and
gravitational terms of Equation~\ref{eq:eom} would lead one to expect: recovering $\tau$
from $\ddot\theta$ alone requires the configuration-dependent inertia and gravity terms
that the kinematic channels do not carry.

\paragraph{Torque contributes a constant increment.} Its effect is close to additive and
independent of which kinematic channels accompany it:
\begin{center}
\begin{tabular}{lcc}
\toprule
Base & \ \ without $\tau$\ \  & \ \ with $\tau$\ \  \\
\midrule
$\sin\theta$                                   & 0.180 & 0.259 \ \ ($+0.079$) \\
$\sin\theta,\cos\theta,\dot\theta,\ddot\theta$ & 0.216 & 0.297 \ \ ($+0.080$) \\
\bottomrule
\end{tabular}
\end{center}
Adding torque is worth $+0.08$ macro mAP whether the loss already scores one kinematic
channel or four. The derivatives behave in the complementary way: they are worth
$+0.058$ when added to the angle alone but only $+0.035$ once torque is present, the
overlap Equation~\ref{eq:eom} implies.

\paragraph{The effect is not explained by supervision breadth.} Transfer is broadly
monotone in head width among the kinematic arms, so a wider prediction target is a
plausible competing account. The dynamic arms refute it directly: the widest kinematic
arm ($32$ dimensions, $0.216$) is beaten by the narrowest arm containing torque ($16$
dimensions, $0.259$), and $\{\sin\theta,\tau\}$ matches the $24$-dimensional
$\{\sin\theta,\cos\theta,\tau\}$ to within the noise floor. What governs transfer is
which physical quantities the loss scores, not how many values it regresses.

We note three limitations. Each ablated arm is a single seed, with the noise floor
estimated from the two seeds of the full objective. All arms are pretrained for 10
epochs, at which point the reconstruction loss is still decreasing, so these are lower
bounds rather than converged values. Finally, the absolute values are specific to the
attentive readout: a mean-pooled probe compresses the same ordering into a narrower
range and roughly triples the seed-to-seed variance.

\subsection{Reading a frozen representation}
\label{sec:pooling}

A linear probe must first reduce the encoder's $T \times D$ token grid to one vector.
This choice is usually treated as an implementation detail; we find it dominates the
measured quality of the representation, and that the conventional default is the worst
of the options we tried.

We compare four readouts over the same frozen backbones, all trained for the same
number of epochs with the same optimizer. \textsc{mean} averages over all $T \cdot D$
tokens. \textsc{mean-max} concatenates the mean with a coordinate-wise maximum,
doubling the readout width without learning anything. \textsc{attentive} scores each
token with a two-layer network and takes a softmax-weighted average. \textsc{factorized}
applies one such scorer across degrees of freedom within each time patch and a second
across time, on the grounds that a token's two coordinates mean different things and a
flat attention has to learn that.

\begin{table}[t]
\centering
\small
\begin{tabular}{lrcccc}
\toprule
Readout & Params & macro mAP & micro mAP & macro F1 & seed spread \\
\midrule
\textsc{mean}        & 15\,420 & 0.289 & 0.501 & 0.131 & 0.024 \\
\textsc{mean-max}    & 30\,780 & 0.260 & 0.458 & 0.057 & 0.011 \\
\textsc{attentive}   & 48\,444 & \textbf{0.341} & 0.529 & 0.196 & 0.002 \\
\textsc{factorized}  & 81\,468 & 0.338 & \textbf{0.529} & \textbf{0.204} & 0.001 \\
\bottomrule
\end{tabular}
\caption{Probe readouts over identical frozen backbones (all-channel objective, 40
pretraining epochs, two seeds). ``Params'' counts every trainable parameter, pooler and
classifier together. ``Seed spread'' is the absolute macro mAP difference between the
two seeds.}
\label{tab:pooling}
\end{table}

\paragraph{The gain is not readout capacity.} \textsc{mean-max} has twice the trainable
parameters of \textsc{mean} and performs worse on every metric, losing more than half of
macro F1 ($0.131 \rightarrow 0.057$). Capacity added along the wrong axis is actively
harmful, so the improvement from attention cannot be attributed to parameter count
alone. \textsc{attentive} improves on \textsc{mean} by $0.052$ macro mAP and $0.065$
macro F1, a relative gain of $49\%$ on the latter.

\paragraph{Factorization is unnecessary.} \textsc{factorized} costs $68\%$ more
parameters than \textsc{attentive} and matches it to within $0.004$ macro mAP, which is
inside the seed spread of either. Attention over the flat grid already recovers whatever
structure the separate spatial and temporal scorers were meant to supply, so we adopt
\textsc{attentive} throughout.

\paragraph{Mean pooling is also unstable.} The two seeds differ by $0.024$ macro mAP
under \textsc{mean} and by $0.002$ under \textsc{attentive}, an order of magnitude. Much
of what would otherwise be reported as run-to-run variance of the pretraining is in fact
variance of the readout. Ablations measured through a mean-pooled probe therefore carry
a noise floor roughly ten times larger than necessary, which matters whenever the effect
under test is comparable in size.

We read this as a statement about what averaging discards. A mean over the whole window
weights every token equally, including the majority that belong to no labelled action;
the localized, short-duration evidence that distinguishes rare classes is diluted in
proportion to window length. That the deficit appears most strongly in macro F1, which
weights rare and frequent classes equally, is consistent with this account. It also
implies that linear-probe numbers are not comparable across papers unless the readout is
specified.
